% =====================================================================
%  BAO CAO SANG KIEN KINH NGHIEM
%  BIEN DICH BANG: XeLaTeX  (khong dung pdflatex vi can font Times New Roman)
%  Lenh: xelatex skkn.tex   (chay 2 lan de co muc luc)
% =====================================================================
\documentclass[a4paper,14pt]{extarticle}

% ----- Font va tieng Viet (XeLaTeX) ----------------------------------
\usepackage{fontspec}
\setmainfont{Times New Roman}
% Neu may bao khong tim thay font, doi thanh:  \setmainfont{TeX Gyre Termes}

% ----- Trinh bay theo quy dinh van ban hanh chinh ---------------------
\usepackage[top=2cm, bottom=2cm, left=3cm, right=2cm]{geometry}
\usepackage{setspace}
\onehalfspacing                      % gian dong 1,5
\setlength{\parindent}{1cm}          % lui dau dong 1cm

% ----- Goi ho tro ----------------------------------------------------
\usepackage{graphicx}
\usepackage{float}
\usepackage{array}
\usepackage{booktabs}
\usepackage{longtable}
\usepackage{multirow}
\usepackage{enumitem}
\usepackage{amsmath, amssymb}
\usepackage{xcolor}
\usepackage{tikz}
\usepackage{titlesec}
\usepackage[hidelinks, unicode]{hyperref}
\usepackage{fancyhdr}
\usepackage[font=small, labelfont=bf, justification=centering]{caption}

\renewcommand{\figurename}{Hình}
\renewcommand{\tablename}{Bảng}
\renewcommand{\contentsname}{MỤC LỤC}
\renewcommand{\refname}{TÀI LIỆU THAM KHẢO}

\pagestyle{fancy}
\fancyhf{}
\fancyfoot[C]{\thepage}
\renewcommand{\headrulewidth}{0pt}

% Danh so muc: I, II, III  -> tieu muc II.1, II.1.1
\renewcommand{\thesection}{\Roman{section}}
\renewcommand{\thesubsection}{\thesection.\arabic{subsection}}
\renewcommand{\thesubsubsection}{\thesubsection.\arabic{subsubsection}}

% Dinh dang tieu de muc
\titleformat{\section}{\normalfont\large\bfseries}{\thesection.}{0.6em}{\MakeUppercase}
\titleformat{\subsection}{\normalfont\normalsize\bfseries}{\thesubsection.}{0.6em}{}
\titleformat{\subsubsection}{\normalfont\normalsize\bfseries\itshape}{\thesubsubsection.}{0.6em}{}

% =====================================================================
%  THONG TIN CAN DIEN  -- sua o day, ca bai tu cap nhat theo
% =====================================================================
\def\SOGD{Sở Giáo dục và Đào tạo \dots\dots}
\def\DONVI{Trường THPT \dots\dots}
\def\TENSANGKIEN{Xây dựng hệ thống Ngân hàng đề Toán ứng dụng trí tuệ nhân tạo
hỗ trợ giáo viên biên soạn đề kiểm tra và tổ chức kiểm tra, đánh giá trực tuyến
ở trường trung học phổ thông}
\def\LINHVUC{Toán học (ứng dụng công nghệ thông tin và chuyển đổi số trong dạy học)}
\def\TACGIA{Mai Hà Lan}
\def\CHUCVU{Giáo viên Toán}
\def\NAMSINH{\dotfill}
\def\TRINHDO{\dotfill}
\def\DIENTHOAI{\dotfill}
\def\EMAIL{lantuan2605@gmail.com}
\def\NAMHOC{2025 -- 2026}
\def\DIADIEM{\dotfill}

% =====================================================================
%  LENH TIEN ICH
% =====================================================================

% --- Bat/tat toan bo ghi chu mau do ---
% Khi in ban nop chinh thuc: doi \hienghichutrue thanh \hienghichufalse
\newif\ifhienghichu
\hienghichutrue

% \ghichu{noi dung nhac viec}
\newcommand{\ghichu}[1]{%
  \ifhienghichu
    \par\smallskip
    \noindent\textcolor{red}{\small\textbf{[GHI CHÚ --- xoá trước khi nộp]}\itshape\ #1}
    \par\smallskip
  \fi}

% \chenanh[do rong]{ten-file.png}{Chu thich hinh}
% Neu file da co trong thu muc hinhanh/ thi chen anh that,
% neu chua co thi hien mot o vuong do bao cho biet cho nay can chup anh gi.
\newcommand{\chenanh}[3][0.85]{%
  \begin{figure}[H]
    \centering
    \IfFileExists{hinhanh/#2}%
      {\includegraphics[width=#1\linewidth]{hinhanh/#2}}%
      {\setlength{\fboxrule}{1.2pt}\setlength{\fboxsep}{12pt}%
       \textcolor{red}{\fbox{\begin{minipage}{#1\linewidth}
         \centering\vspace{1.2cm}
         \textbf{CHỖ CHÈN ẢNH}\\[6pt]
         \texttt{hinhanh/#2}\\[10pt]
         \itshape\small Cần chụp: #3
         \vspace{1.2cm}
       \end{minipage}}}}%
    \caption{#3}
  \end{figure}}

% =====================================================================
\begin{document}

% ---------------------------------------------------------------
%  BIA SANG KIEN  (Mau 2 theo huong dan cua So GD&DT)
% ---------------------------------------------------------------
\thispagestyle{empty}
\begin{center}
\setstretch{1.2}

{\large\bfseries \MakeUppercase{\SOGD}}\\[4pt]
{\large\bfseries \MakeUppercase{\DONVI}}\\[2.5cm]

\includegraphics[width=2.5cm]{example-image-a}%
\ghichu{Dòng trên đang dùng ảnh mẫu của LaTeX. Nếu trường có logo, lưu logo vào
\texttt{hinhanh/logo.png} rồi đổi \texttt{example-image-a} thành \texttt{hinhanh/logo.png}.
Nếu không cần logo thì xoá hẳn dòng \texttt{includegraphics} này.}

\vspace{1.5cm}

{\Large\bfseries BÁO CÁO SÁNG KIẾN}\\[1.2cm]

{\LARGE\bfseries \TENSANGKIEN}\\[2cm]

\end{center}

\begin{flushleft}
\hspace{5cm}\begin{minipage}{0.55\linewidth}
\setstretch{1.5}
\textbf{Lĩnh vực:} \LINHVUC

\textbf{Tác giả:} \TACGIA

\textbf{Chức vụ:} \CHUCVU

\textbf{Đơn vị công tác:} \DONVI

\textbf{Năm học:} \NAMHOC
\end{minipage}
\end{flushleft}

\vfill
\begin{center}
{\bfseries \DIADIEM, tháng \dots\ năm 2026}
\end{center}

\clearpage

% ---------------------------------------------------------------
%  DON YEU CAU CONG NHAN SANG KIEN  (Mau 3)
% ---------------------------------------------------------------
\thispagestyle{empty}
\begin{center}
{\bfseries CỘNG HÒA XÃ HỘI CHỦ NGHĨA VIỆT NAM}\\
{\bfseries Độc lập -- Tự do -- Hạnh phúc}\\[4pt]
\rule{5cm}{0.5pt}\\[1cm]
{\Large\bfseries ĐƠN YÊU CẦU CÔNG NHẬN SÁNG KIẾN}\\[1cm]
\end{center}

\noindent\textbf{Kính gửi:} Hội đồng Sáng kiến \dotfill
\ghichu{Điền đúng tên hội đồng ghi trong công văn hướng dẫn của Sở
(thường là ``Hội đồng Sáng kiến ngành Giáo dục và Đào tạo tỉnh \dots'').}

\medskip
\noindent Tôi ghi tên dưới đây:

\begin{center}
\begin{tabular}{|c|p{2.6cm}|p{2cm}|p{2.4cm}|p{1.8cm}|p{2cm}|p{1.5cm}|}
\hline
\textbf{TT} & \textbf{Họ và tên} & \textbf{Ngày sinh} & \textbf{Nơi công tác} &
\textbf{Chức danh} & \textbf{Trình độ chuyên môn} & \textbf{Tỷ lệ đóng góp} \\ \hline
1 & \TACGIA & \NAMSINH & \DONVI & \CHUCVU & \TRINHDO & 100\% \\ \hline
\end{tabular}
\end{center}

\ghichu{Điền ngày sinh, trình độ chuyên môn (Cử nhân/Thạc sĩ Toán học) vào file
\texttt{skkn.tex} ở phần ``THÔNG TIN CẦN ĐIỀN'', cả bài sẽ tự cập nhật.}

\bigskip
\noindent\textbf{1. Là tác giả đề nghị xét công nhận sáng kiến:} \TENSANGKIEN

\noindent\textbf{2. Chủ đầu tư tạo ra sáng kiến:} \TACGIA\ (tác giả tự đầu tư
kinh phí thuê máy chủ, tên miền và tự thực hiện toàn bộ phần lập trình).

\noindent\textbf{3. Lĩnh vực áp dụng sáng kiến:} \LINHVUC. Sáng kiến được áp
dụng trong khâu biên soạn đề kiểm tra và tổ chức kiểm tra, đánh giá môn Toán
cấp trung học phổ thông, phục vụ đồng thời giáo viên (lấy đề, giao bài, theo
dõi kết quả lớp) và học sinh (tự luyện tập, làm bài trực tuyến).

\noindent\textbf{4. Ngày sáng kiến được áp dụng lần đầu hoặc áp dụng thử:}
\dotfill
\ghichu{Ghi ngày chị bắt đầu cho học sinh làm bài trên hệ thống. Nếu không nhớ
chính xác, lấy mốc đầu học kỳ (ví dụ 05/9/2025).}

\bigskip
\noindent\textbf{5. Mô tả bản chất của sáng kiến}

\subsection*{5.1. Tình trạng giải pháp đã biết}

Việc biên soạn đề kiểm tra môn Toán ở trường phổ thông hiện nay chủ yếu được
thực hiện thủ công: giáo viên sưu tầm câu hỏi từ nhiều nguồn rời rạc, gõ lại
bằng Word hoặc \LaTeX, tự cân đối ma trận rồi tự tạo các mã đề bằng cách đảo
thứ tự câu và phương án. Cách làm này bộc lộ bốn hạn chế: mất nhiều thời gian;
khó bảo đảm đúng ma trận theo bốn mức độ nhận thức; các mã đề chỉ khác nhau về
thứ tự nên thực chất vẫn là một đề; và kho câu hỏi không được mã hoá thống nhất
nên gần như không tái sử dụng được ở các năm học sau.

Về phía học sinh, nhu cầu tự luyện tập không được đáp ứng đúng cách: các em
chỉ tìm được đề trôi nổi trên mạng, thường lệch phạm vi chương trình hoặc sai
lời giải; làm xong phải chờ đến tiết sau mới biết đúng sai nên phản hồi mất
tính thời sự; và không có cách nào luyện lại nhiều lần cùng một dạng toán với
số liệu khác nhau.

Một số phần mềm tạo đề thương mại đã có trên thị trường nhưng hoặc thu phí,
hoặc chỉ hoán vị câu hỏi trong một kho có sẵn của nhà cung cấp, không cho phép
giáo viên đưa ngân hàng câu hỏi của chính mình vào, và chưa cập nhật kịp cấu
trúc định dạng đề mới gồm bốn dạng câu hỏi (trắc nghiệm nhiều lựa chọn,
đúng/sai, trả lời ngắn, tự luận) theo Quyết định số 764/QĐ-BGDĐT ngày
08/3/2024 của Bộ Giáo dục và Đào tạo. (Phân tích chi tiết ở Mục II.2.)

\subsection*{5.2. Nội dung giải pháp đề nghị công nhận là sáng kiến}

\textbf{a) Mục đích của giải pháp}

Xây dựng một hệ thống ngân hàng đề môn Toán chạy trên nền web dùng chung cho
cả giáo viên và học sinh. Về phía học sinh: chủ động tạo đề luyện tập đúng phạm
vi mình cần bằng một câu yêu cầu tiếng Việt, làm bài trực tuyến và nhận điểm
cùng đáp án ngay. Về phía giáo viên: có đề kiểm tra đúng ma trận chỉ sau vài
phút, mỗi lần tạo cho ra một đề thực sự khác nhau về số liệu, giao đề và trả
điểm về lớp học Google Classroom, đồng thời theo dõi được điểm trung bình và
những bài học sinh hay sai.

\textbf{b) Tính mới của giải pháp}

\begin{enumerate}[label=(\arabic*), leftmargin=1.2cm, itemsep=2pt]
  \item Câu hỏi không được lưu ở dạng văn bản tĩnh mà được lưu dưới dạng
        \emph{hàm sinh} viết bằng Python. Mỗi lần gọi, hàm sinh ra một câu hỏi
        mới với bộ số liệu mới và lời giải tương ứng, nên từ một câu gốc có thể
        tạo ra hàng trăm biến thể cùng dạng, cùng độ khó.
  \item Toàn bộ kho câu hỏi được mã hoá theo một chuẩn định danh thống nhất
        gắn với lớp -- chương -- bài -- mức độ -- dạng câu, nhờ đó máy có thể
        tự chọn câu theo ma trận mà không cần giáo viên chỉ định từng câu.
  \item Cùng một ngân hàng câu hỏi phục vụ hai đối tượng: dữ liệu học sinh
        làm bài quay trở lại thành thông tin thống kê cho giáo viên, tạo thành
        một vòng khép kín giữa dạy và học.
  \item Thuật toán phân bổ câu hỏi bám theo số tiết của từng bài trong phân
        phối chương trình, đồng thời xử lý riêng dạng câu đúng/sai theo quy
        ước một câu đúng/sai tương đương bốn ý ở bốn mức độ khác nhau.
  \item Trí tuệ nhân tạo chỉ được dùng ở khâu hiểu yêu cầu bằng ngôn ngữ tự
        nhiên của giáo viên; toàn bộ khâu chọn câu, tính điểm và sinh đề do mã
        nguồn đảm nhiệm. Nguyên tắc ``ưu tiên mã nguồn hơn trí tuệ nhân tạo''
        này giúp kết quả ổn định, kiểm chứng được và không phụ thuộc vào việc
        mô hình có trả lời đúng hay không.
  \item Hệ thống khép kín một vòng: từ yêu cầu của giáo viên đến đề PDF, lời
        giải PDF, bài làm trực tuyến, điểm số và sổ điểm trên Google Classroom.
\end{enumerate}

\textbf{c) Cách thức thực hiện}

Trình bày chi tiết tại Mục II.3 của báo cáo, gồm sáu bước: chuẩn hoá mã câu
hỏi; viết hàm sinh câu hỏi bằng Python; xây dựng bảng ma trận đề; lập trình
thuật toán chọn câu; sinh đề và lời giải bằng \LaTeX; tổ chức làm bài, chấm
điểm và trả kết quả về Google Classroom.

\subsection*{5.3. Khả năng áp dụng của giải pháp}

Sáng kiến đã được áp dụng tại \DONVI\ cho cả giáo viên tổ Toán và học sinh,
và có thể áp dụng cho mọi tổ Toán ở trường trung học cơ sở, trung học phổ
thông. Do phần khung hệ thống tách rời
khỏi nội dung câu hỏi, các môn học khác có ngân hàng câu hỏi dạng tính toán
(Vật lí, Hoá học, Sinh học) có thể dùng lại toàn bộ phần mềm, chỉ cần bổ sung
hàm sinh câu hỏi của môn mình.

\subsection*{5.4. Hiệu quả, lợi ích thu được}

\ghichu{Mục này em sẽ dựng bảng số liệu ở phần sau; chị điền số thật rồi chép
tóm tắt hai đến ba dòng vào đây. Không nên ghi số ước lượng trong đơn.}

Trình bày chi tiết tại Mục II.5 của báo cáo.

\subsection*{5.5. Tài liệu kèm theo}

Báo cáo sáng kiến; phụ lục hình ảnh sản phẩm; địa chỉ truy cập hệ thống;
tài khoản dùng thử dành cho hội đồng chấm.

\bigskip
\noindent\textbf{6. Những thông tin cần được bảo mật:} Không có.

\noindent\textbf{7. Các điều kiện cần thiết để áp dụng sáng kiến:} Máy tính
hoặc điện thoại có kết nối Internet; giáo viên biết sử dụng máy tính ở mức cơ
bản; nhà trường có lớp học Google Classroom (không bắt buộc).

\bigskip
\noindent Tôi xin cam đoan mọi thông tin nêu trong đơn là trung thực, đúng sự
thật và hoàn toàn chịu trách nhiệm trước pháp luật.

\vspace{1cm}
\begin{flushright}
\begin{minipage}{6cm}
\centering
\textit{\DIADIEM, ngày \dots\ tháng \dots\ năm 2026}\\[2pt]
\textbf{Người nộp đơn}\\
\textit{(Ký và ghi rõ họ tên)}\\[2.5cm]
\TACGIA
\end{minipage}
\end{flushright}

\clearpage


\clearpage
\tableofcontents
\clearpage

% ---------------------------------------------------------------
%  PHAN I -- MO DAU
% ---------------------------------------------------------------
\section{MỞ ĐẦU}

\subsection{Lí do chọn đề tài}

Kiểm tra, đánh giá là khâu then chốt quyết định chất lượng dạy học. Chương
trình Giáo dục phổ thông 2018 ban hành kèm theo Thông tư số 32/2018/TT-BGDĐT
chuyển trọng tâm đánh giá từ ghi nhớ kiến thức sang đánh giá năng lực; Thông tư
số 22/2021/TT-BGDĐT quy định việc đánh giá học sinh trung học cơ sở và trung
học phổ thông phải bảo đảm tính chính xác, công bằng, vì sự tiến bộ của người
học. Đặc biệt, Quyết định số 764/QĐ-BGDĐT ngày 08/3/2024 của Bộ Giáo dục và Đào
tạo công bố cấu trúc định dạng đề thi tốt nghiệp trung học phổ thông từ năm
2025 với ba dạng câu hỏi trắc nghiệm mới --- nhiều lựa chọn, đúng/sai và trả
lời ngắn --- đã đặt ra yêu cầu mới cho mỗi giáo viên Toán: đề kiểm tra thường
xuyên, định kì ở trường phải sớm làm quen học sinh với cấu trúc đó.

Thực tế tại trường tôi, việc đáp ứng yêu cầu này gặp ba khó khăn lớn. Thứ nhất,
soạn một đề kiểm tra đúng ma trận bốn mức độ với đủ bốn dạng câu hỏi mất từ bốn
đến sáu giờ làm việc, trong khi một giáo viên thường phải soạn nhiều đề cho
nhiều lớp trong cùng một đợt kiểm tra. Thứ hai, các mã đề được tạo bằng cách
đảo thứ tự câu hỏi và phương án nên thực chất vẫn là một đề, học sinh ngồi gần
nhau vẫn có thể trao đổi đáp án. Thứ ba, kho câu hỏi tích luỹ qua nhiều năm nằm
rải rác trong hàng trăm tệp Word và \LaTeX\ không được đặt tên thống nhất, đến
năm sau gần như phải soạn lại từ đầu.

Trong khi đó, Nghị quyết số 57-NQ/TW ngày 22/12/2024 của Bộ Chính trị về đột
phá phát triển khoa học, công nghệ, đổi mới sáng tạo và chuyển đổi số quốc gia
xác định chuyển đổi số là phương thức sản xuất mới, yêu cầu ứng dụng mạnh mẽ
trí tuệ nhân tạo vào mọi lĩnh vực, trong đó có giáo dục. Các mô hình ngôn ngữ
lớn hiện nay đã đủ khả năng hiểu yêu cầu bằng tiếng Việt của giáo viên; tuy
nhiên nếu giao hẳn việc ra đề cho trí tuệ nhân tạo thì lời giải thường sai và
không kiểm chứng được, không thể dùng làm đề kiểm tra chính thức.

Từ ba khó khăn thực tế và một cơ hội công nghệ nêu trên, tôi chọn nghiên cứu và
thực hiện sáng kiến \textit{``\TENSANGKIEN''}, với hướng đi khác các phần mềm
hiện có: trí tuệ nhân tạo chỉ làm nhiệm vụ hiểu yêu cầu của giáo viên, còn việc
sinh câu hỏi, chọn câu theo ma trận và tính điểm đều do mã nguồn đảm nhiệm để
bảo đảm tính chính xác tuyệt đối.

\ghichu{Nếu trường chị có đặc thù riêng (ví dụ lớp đông, thiếu phòng máy, học
sinh vùng khó) thì thêm một hai câu vào đoạn ``Thực tế tại trường tôi'' --- hội
đồng chấm rất thích chi tiết cụ thể của đơn vị.}

\subsection{Mục tiêu của sáng kiến}

\textbf{Mục tiêu chung:} Xây dựng và đưa vào sử dụng một hệ thống ngân hàng đề
môn Toán trên nền web, giúp giáo viên rút ngắn thời gian biên soạn đề kiểm tra,
bảo đảm đề đúng ma trận và mỗi học sinh làm một đề riêng, đồng thời hỗ trợ tổ
chức kiểm tra và chấm điểm trực tuyến.

\textbf{Mục tiêu cụ thể:}

\begin{enumerate}[label=(\arabic*), leftmargin=1.2cm, itemsep=2pt]
  \item Xây dựng chuẩn mã hoá định danh câu hỏi thống nhất theo lớp, chương,
        bài, mức độ nhận thức và dạng câu hỏi, làm cơ sở để máy tự chọn câu.
  \item Xây dựng thư viện hàm sinh câu hỏi bằng Python cho phép tạo ra vô số
        biến thể của cùng một dạng toán, kèm lời giải chi tiết chính xác.
  \item Lập trình thuật toán phân bổ câu hỏi theo ma trận đề, bám số tiết của
        từng bài trong phân phối chương trình.
  \item Tự động sinh đề và lời giải dưới dạng tệp PDF trình bày bằng \LaTeX,
        đạt chuẩn văn bản đề thi.
  \item Triển khai chức năng làm bài, chấm điểm tự động theo thang điểm 10 và
        trả đề, đáp án, điểm số về lớp học Google Classroom của học sinh.
  \item Đánh giá hiệu quả của hệ thống qua thời gian soạn đề của giáo viên và
        kết quả học tập của học sinh.
\end{enumerate}

\subsection{Đối tượng, phạm vi và thời gian thực hiện}

\textbf{Đối tượng nghiên cứu:} Quy trình biên soạn đề kiểm tra và tổ chức kiểm
tra, đánh giá môn Toán ở trường trung học phổ thông; giải pháp tự động hoá quy
trình đó bằng ngân hàng câu hỏi tham số và trí tuệ nhân tạo.

\textbf{Đối tượng áp dụng:} Giáo viên tổ Toán và học sinh các lớp \dotfill\ của
\DONVI.
\ghichu{Điền tên lớp thực nghiệm và lớp đối chứng, ví dụ ``lớp 10A1, 10A2''.}

\textbf{Phạm vi nội dung:} Sáng kiến được triển khai trước hết với chương trình
Toán lớp 10 (bộ sách Kết nối tri thức với cuộc sống), chương \dotfill; phần
khung hệ thống áp dụng được cho mọi khối lớp.
\ghichu{Ghi rõ chương đã hoàn thành ngân hàng câu hỏi, ví dụ ``Chương I. Mệnh
đề và tập hợp''. Không nên ghi rộng hơn phần thực sự đã làm --- hội đồng sẽ
kiểm tra.}

\textbf{Thời gian thực hiện:} Từ tháng \dots\ năm 2025 đến tháng \dots\ năm
2026, gồm ba giai đoạn: nghiên cứu và thiết kế hệ thống; lập trình và xây dựng
ngân hàng câu hỏi; thực nghiệm sư phạm và hoàn thiện.

\subsection{Phương pháp nghiên cứu}

\begin{itemize}[leftmargin=1.2cm, itemsep=2pt]
  \item \textbf{Phương pháp nghiên cứu lí luận:} nghiên cứu Chương trình Giáo
        dục phổ thông 2018, các văn bản chỉ đạo về kiểm tra đánh giá và cấu
        trúc định dạng đề thi, tài liệu về đo lường và đánh giá trong giáo dục.
  \item \textbf{Phương pháp điều tra, khảo sát:} khảo sát giáo viên tổ Toán về
        thời gian và khó khăn khi biên soạn đề; khảo sát học sinh về mức độ
        hứng thú và tính công bằng của hình thức kiểm tra.
  \item \textbf{Phương pháp thực nghiệm sư phạm:} tổ chức dạy học và kiểm tra
        có đối chứng giữa lớp sử dụng hệ thống và lớp không sử dụng.
  \item \textbf{Phương pháp thống kê toán học:} xử lí số liệu điểm kiểm tra để
        so sánh kết quả hai nhóm lớp.
  \item \textbf{Phương pháp chuyên gia:} xin ý kiến tổ chuyên môn và ban giám
        hiệu về chất lượng đề do hệ thống sinh ra.
\end{itemize}

\subsection{Điểm mới của sáng kiến}

So với các giải pháp đã có, sáng kiến có bốn điểm mới:

\begin{enumerate}[label=(\arabic*), leftmargin=1.2cm, itemsep=2pt]
  \item \textbf{Câu hỏi tham số thay cho câu hỏi tĩnh.} Mỗi câu hỏi trong ngân
        hàng là một chương trình sinh câu hỏi, không phải một đoạn văn bản cố
        định; nhờ đó mỗi học sinh nhận một đề khác nhau về số liệu nhưng tương
        đương về dạng toán và độ khó.
  \item \textbf{Phân vai rõ ràng giữa trí tuệ nhân tạo và mã nguồn.} Trí tuệ
        nhân tạo chỉ đọc hiểu yêu cầu bằng tiếng Việt của giáo viên; việc chọn
        câu, tính điểm và tạo đề hoàn toàn do mã nguồn thực hiện nên kết quả
        chính xác và lặp lại được.
  \item \textbf{Ma trận đề đặt trong bảng cấu hình, không gài cứng trong mã.}
        Khi Bộ hoặc Sở thay đổi cấu trúc đề, giáo viên chỉ sửa số liệu trong
        bảng mà không cần lập trình lại.
  \item \textbf{Khép kín một vòng dạy học.} Từ một câu yêu cầu bằng tiếng Việt
        đến đề PDF, lời giải, bài làm trực tuyến, điểm số và sổ điểm Google
        Classroom đều diễn ra trong cùng một hệ thống.
\end{enumerate}

\chenanh[0.9]{01-giao-dien-tong-quan.png}{Giao diện trang chủ hệ thống Ngân hàng đề Toán tại địa chỉ nganhangdechv.tech}

\ghichu{Ảnh 01: chụp toàn màn hình trang chủ sau khi đăng nhập, thấy rõ ô chat
và tên hệ thống. Lưu thành \texttt{skkn/hinhanh/01-giao-dien-tong-quan.png}.
Khi đã có file, ô đỏ sẽ tự biến mất và ảnh hiện ra --- không phải sửa gì trong
bài.}

\clearpage

% ---------------------------------------------------------------
%  PHAN II -- NOI DUNG
%  Muc II.1 -- Co so li luan va co so phap li
% ---------------------------------------------------------------
\section{NỘI DUNG SÁNG KIẾN}

\subsection{Cơ sở lí luận và cơ sở pháp lí}

\subsubsection{Cơ sở pháp lí}

Sáng kiến được xây dựng trên các căn cứ pháp lí sau:

\begin{itemize}[leftmargin=1.2cm, itemsep=3pt]
  \item \textbf{Thông tư số 32/2018/TT-BGDĐT} ngày 26/12/2018 ban hành Chương
        trình Giáo dục phổ thông. Chương trình môn Toán xác định mục tiêu hình
        thành và phát triển năm thành tố của năng lực toán học, trong đó có
        năng lực giải quyết vấn đề toán học và năng lực sử dụng công cụ,
        phương tiện học toán. Đây là căn cứ để thiết kế câu hỏi theo mức độ
        năng lực chứ không theo mức độ ghi nhớ.

  \item \textbf{Thông tư số 22/2021/TT-BGDĐT} ngày 20/7/2021 quy định về đánh
        giá học sinh trung học cơ sở và trung học phổ thông. Thông tư yêu cầu
        việc đánh giá phải \textit{``vì sự tiến bộ của học sinh''}, coi trọng
        động viên, khuyến khích sự cố gắng của người học, bảo đảm kịp thời,
        công bằng, khách quan. Hai yêu cầu \textit{kịp thời} và \textit{công
        bằng} chính là hai vấn đề mà sáng kiến này hướng tới giải quyết bằng
        công nghệ: chấm ngay khi học sinh nộp bài, và mỗi học sinh một đề
        riêng.

  \item \textbf{Quyết định số 764/QĐ-BGDĐT} ngày 08/3/2024 của Bộ Giáo dục và
        Đào tạo về việc phê duyệt cấu trúc định dạng đề thi tốt nghiệp trung
        học phổ thông từ năm 2025, với ba dạng câu hỏi trắc nghiệm: nhiều lựa
        chọn, đúng/sai và trả lời ngắn. Đây là căn cứ để hệ thống sinh đủ bốn
        dạng câu hỏi (ba dạng trắc nghiệm nêu trên và tự luận).

  \item \textbf{Quyết định số 131/QĐ-TTg} ngày 25/01/2022 của Thủ tướng Chính
        phủ phê duyệt Đề án ``Tăng cường ứng dụng công nghệ thông tin và
        chuyển đổi số trong giáo dục và đào tạo''.

  \item \textbf{Nghị quyết số 57-NQ/TW} ngày 22/12/2024 của Bộ Chính trị về
        đột phá phát triển khoa học, công nghệ, đổi mới sáng tạo và chuyển đổi
        số quốc gia, trong đó xác định phải đẩy mạnh ứng dụng trí tuệ nhân tạo
        và dữ liệu lớn trong các lĩnh vực, đưa giáo dục trở thành một trong
        những lĩnh vực được ưu tiên chuyển đổi số.

  \item \textbf{Nghị định số 13/2012/NĐ-CP} ngày 02/3/2012 của Chính phủ ban
        hành Điều lệ Sáng kiến và các văn bản hướng dẫn hiện hành của Sở Giáo
        dục và Đào tạo về công tác sáng kiến năm học \NAMHOC.
\end{itemize}

\ghichu{Thay dòng cuối bằng đúng số hiệu công văn hướng dẫn sáng kiến năm nay
của Sở mình --- hội đồng chấm thường đối chiếu ngay dòng này.}

\subsubsection{Cơ sở lí luận}

\textbf{a) Đánh giá vì sự tiến bộ của người học đòi hỏi phản hồi tức thời}

Lí luận dạy học hiện đại phân biệt \textit{đánh giá kết quả học tập}
(\textit{assessment of learning}) với \textit{đánh giá vì học tập}
(\textit{assessment for learning}). Ở dạng thứ hai, giá trị của bài kiểm tra
không nằm ở điểm số mà nằm ở thông tin phản hồi mà học sinh nhận được và ở
việc học sinh được luyện lại ngay phần mình còn yếu. Muốn vậy phải đáp ứng hai
điều kiện: phản hồi phải đến \textit{ngay} sau khi làm bài, và học sinh phải
có \textit{đủ đề} để luyện lại nhiều lần cùng một dạng toán. Cách làm thủ công
không đáp ứng được cả hai điều kiện này --- giáo viên không thể chấm ngay cho
mấy trăm bài, cũng không thể soạn cho mỗi em một đề luyện tập riêng.

\textbf{b) Câu hỏi tham số hoá và tính tương đương của đề}

Trong đo lường giáo dục, hai câu hỏi được coi là \textit{tương đương} khi
chúng cùng đo một chỉ báo năng lực, cùng dạng thao tác tư duy và cùng mức độ
nhận thức, dù số liệu cụ thể khác nhau. Từ nhận xét đó hình thành kĩ thuật
\textit{câu hỏi tham số hoá} (\textit{parameterized item} hay
\textit{algorithmic item}): thay vì lưu một câu hỏi cố định, người ta lưu quy
tắc sinh ra câu hỏi cùng với quy tắc sinh đáp án và lời giải. Mỗi lần gọi, quy
tắc cho ra một câu hỏi mới với bộ số liệu mới nhưng vẫn giữ nguyên dạng toán
và độ khó.

Kĩ thuật này giải quyết triệt để vấn đề công bằng trong kiểm tra: các mã đề
không còn là một đề hoán vị thứ tự mà là những đề thực sự khác nhau về số
liệu. Đồng thời nó biến kho câu hỏi từ tài nguyên \textit{hữu hạn} thành tài
nguyên \textit{không cạn}: từ một câu gốc có thể sinh ra hàng trăm câu luyện
tập cho học sinh.

Điều kiện để áp dụng được kĩ thuật này là bộ tham số phải được ràng buộc chặt
để mọi biến thể đều có nghiệm đẹp, đúng phạm vi chương trình và không rơi vào
trường hợp suy biến --- đây chính là phần công sức chuyên môn của giáo viên bộ
môn, máy không làm thay được.

\textbf{c) Ma trận đề và bốn mức độ nhận thức}

Đề kiểm tra chỉ có giá trị đo lường khi được xây dựng từ một ma trận xác định
rõ tỉ lệ câu hỏi theo từng đơn vị kiến thức và từng mức độ nhận thức: nhận
biết, thông hiểu, vận dụng, vận dụng cao. Về nguyên tắc, trọng số của một đơn
vị kiến thức trong đề phải tương ứng với thời lượng dạy học của đơn vị đó
trong phân phối chương trình. Nguyên tắc này thường bị bỏ qua khi soạn đề thủ
công vì phải tính toán thủ công quá nhiều; ngược lại, đó lại là việc máy tính
làm rất tốt và tuyệt đối chính xác.

\textbf{d) Trí tuệ nhân tạo là trợ lí, không phải người thay thế giáo viên}

Các mô hình ngôn ngữ lớn hiện nay rất mạnh ở việc hiểu ngôn ngữ tự nhiên,
nhưng bản chất là mô hình dự đoán chuỗi kí tự có xác suất cao nhất, không phải
hệ suy luận toán học. Khi được yêu cầu tự ra đề và giải, mô hình có thể cho
lời giải trôi chảy nhưng sai ở bước tính toán, và lần chạy sau lại cho kết quả
khác lần trước. Một đề kiểm tra lấy điểm chính thức không chấp nhận được rủi
ro đó.

Vì vậy, sáng kiến này xác định rõ ranh giới phân vai: trí tuệ nhân tạo chỉ làm
những việc \textit{cần hiểu ngôn ngữ tự nhiên hoặc cần văn phong tự nhiên};
mọi việc có quy tắc rõ ràng --- tra bảng phân phối chương trình, dựng ma trận,
chọn mã câu hỏi, sinh số liệu, tính điểm, so khớp đáp án --- đều do mã nguồn
thực hiện. Còn quyền quyết định cuối cùng về nội dung và mức độ đề vẫn thuộc
về giáo viên. Nguyên tắc này được ghi thành quy tắc phát triển bắt buộc của hệ
thống: \textit{``Ưu tiên mã nguồn hơn trí tuệ nhân tạo; không để nghiệp vụ nằm
trong câu lệnh gửi cho mô hình.''}

\begin{figure}[H]
\centering
\begin{tikzpicture}[
  font=\small,
  box/.style={draw, rounded corners=3pt, align=center, minimum height=1.1cm,
              text width=3.3cm, thick},
  ai/.style={box, fill=blue!8},
  code/.style={box, fill=green!8},
  gv/.style={box, fill=orange!12},
  mui/.style={-latex, thick}
]
\node[gv] (gv) at (0,0) {\textbf{Giáo viên}\\[2pt]\scriptsize quyết định nội dung, mức độ, duyệt đề};
\node[ai] (ai) at (5,0) {\textbf{Trí tuệ nhân tạo}\\[2pt]\scriptsize hiểu yêu cầu bằng tiếng Việt};
\node[code] (code) at (10,0) {\textbf{Mã nguồn}\\[2pt]\scriptsize chọn câu, sinh số liệu, tính điểm};
\draw[mui] (gv) -- (ai);
\draw[mui] (ai) -- (code);
\node[align=center, font=\scriptsize\itshape] at (5,-1.6)
  {chỉ khâu ngôn ngữ};
\node[align=center, font=\scriptsize\itshape] at (10,-1.6)
  {mọi khâu có quy tắc};
\end{tikzpicture}
\caption{Phân vai giữa giáo viên, trí tuệ nhân tạo và mã nguồn trong hệ thống}
\end{figure}

\textbf{e) Một hệ thống phục vụ đồng thời hai đối tượng}

Từ bốn cơ sở trên, hệ thống được thiết kế để phục vụ đồng thời giáo viên và
học sinh, với hai vai trò khác nhau nhưng dùng chung một ngân hàng câu hỏi:

\begin{itemize}[leftmargin=1.2cm, itemsep=3pt]
  \item \textbf{Về phía học sinh:} chủ động tạo đề luyện tập đúng phạm vi mình
        cần bằng một câu yêu cầu tiếng Việt, làm bài trực tuyến và nhận điểm
        cùng đáp án ngay. Đây là điều kiện để đánh giá thường xuyên thực sự
        diễn ra \textit{vì sự tiến bộ} của các em.
  \item \textbf{Về phía giáo viên:} có nguồn đề đúng ma trận để dùng cho kiểm
        tra trên lớp, giao đề và trả điểm qua Google Classroom, đồng thời theo
        dõi được điểm trung bình và những bài học sinh hay sai để điều chỉnh
        việc dạy.
\end{itemize}

Cách phân vai này cũng phù hợp với triết lí phát triển đã đặt ra ngay từ đầu
của hệ thống: \textit{``Không dùng trí tuệ nhân tạo để thay thế giáo viên.
Trí tuệ nhân tạo đóng vai trò trợ lí. Giáo viên vẫn là người quyết định.''}

\clearpage

% % ---------------------------------------------------------------
%  Muc II.2 -- Thuc trang (tinh trang giai phap da biet)
% ---------------------------------------------------------------
\subsection{Thực trạng vấn đề --- tình trạng các giải pháp đã biết}

\subsubsection{Thực trạng biên soạn đề kiểm tra tại đơn vị}

Để có căn cứ khách quan, tôi đã khảo sát giáo viên tổ Toán của nhà trường
trước khi triển khai sáng kiến.

\begin{center}
\begin{tabular}{|c|p{8.2cm}|c|c|}
\hline
\textbf{TT} & \textbf{Nội dung khảo sát} & \textbf{Số GV} & \textbf{Tỉ lệ} \\ \hline
1 & Thời gian soạn một đề kiểm tra định kì trên 4 giờ & & \\ \hline
2 & Cho rằng khó bảo đảm đúng tỉ lệ bốn mức độ nhận thức & & \\ \hline
3 & Tạo mã đề bằng cách đảo thứ tự câu hỏi và phương án & & \\ \hline
4 & Không tái sử dụng được kho câu hỏi của năm trước & & \\ \hline
5 & Chưa quen soạn đủ bốn dạng câu hỏi theo cấu trúc mới & & \\ \hline
\end{tabular}
\end{center}
\textit{Bảng 1. Kết quả khảo sát giáo viên tổ Toán (tổng số: \dots\ giáo viên)}

\ghichu{Chị phát phiếu cho tổ Toán rồi điền hai cột cuối. Nếu tổ ít người, ghi
thẳng số lượng (ví dụ 5/7) --- hội đồng không đòi mẫu lớn, chỉ đòi có số thật.
Nếu chưa kịp khảo sát, có thể bỏ bảng này và viết thành đoạn văn nhận định,
nhưng có bảng thì thuyết phục hơn nhiều.}

Qua trao đổi trực tiếp, các giáo viên đều cho rằng khâu tốn công nhất không
phải là nghĩ ra câu hỏi mà là \textit{cân đối ma trận} và \textit{tạo các mã đề
khác nhau}. Một đề kiểm tra định kì theo cấu trúc mới gồm 12 câu trắc nghiệm
nhiều lựa chọn, 2 câu đúng/sai (tương đương 8 ý), 4 câu trả lời ngắn và 3 câu
tự luận; để có bốn mã đề thực sự khác nhau thì phải soạn tới bốn đề độc lập ---
một khối lượng công việc không khả thi trong điều kiện dạy học bình thường.

\subsubsection{Thực trạng tự luyện tập của học sinh}

\begin{center}
\begin{tabular}{|c|p{8.2cm}|c|c|}
\hline
\textbf{TT} & \textbf{Nội dung khảo sát} & \textbf{Số HS} & \textbf{Tỉ lệ} \\ \hline
1 & Có nhu cầu tự luyện thêm bài tập sau mỗi bài học & & \\ \hline
2 & Tìm đề luyện tập trên mạng, không rõ có đúng chương trình không & & \\ \hline
3 & Gặp lời giải sai hoặc không hiểu được khi tự tra trên mạng & & \\ \hline
4 & Muốn biết ngay đúng/sai sau khi làm bài & & \\ \hline
5 & Không biết mình đang yếu ở bài nào cụ thể & & \\ \hline
\end{tabular}
\end{center}
\textit{Bảng 2. Kết quả khảo sát học sinh trước khi áp dụng (tổng số: \dots\ học sinh)}

\ghichu{Khảo sát nhanh bằng Google Biểu mẫu trong 5 phút đầu tiết là đủ. Nhớ
giữ lại ảnh chụp biểu mẫu và bảng kết quả để đưa vào phụ lục --- đó là bằng
chứng tốt.}

Điểm đáng chú ý nhất là học sinh không thiếu \textit{ý chí} luyện tập mà thiếu
\textit{nguồn đề đúng tầm}: đề trên mạng hoặc quá khó, hoặc thuộc chương trình
cũ, hoặc có lời giải sai; và dù có làm thì cũng phải chờ giáo viên chữa mới
biết mình đúng hay sai.

\subsubsection{Phân tích các giải pháp đã có}

Trên thực tế hiện đã có ba nhóm giải pháp, mỗi nhóm giải quyết được một phần và
để lại một hạn chế cốt lõi.

\textbf{Nhóm 1 --- Ngân hàng đề tĩnh do giáo viên tự tập hợp.} Giáo viên sưu
tầm câu hỏi từ nhiều nguồn, gõ vào Word hoặc \LaTeX\ rồi cắt ghép thành đề.
Ưu điểm là nội dung do chính giáo viên kiểm soát. Hạn chế cốt lõi: kho câu hỏi
là \textit{hữu hạn và tĩnh} --- dùng một lần là học sinh biết đề, muốn có đề
mới phải ngồi thay số bằng tay từng câu một, và mỗi lần thay số lại phải giải
lại để kiểm tra đáp án.

\textbf{Nhóm 2 --- Phần mềm trộn đề và các website ngân hàng đề thương mại.}
Các phần mềm này lấy câu hỏi từ một kho có sẵn của nhà cung cấp rồi hoán vị thứ
tự câu và thứ tự phương án. Hạn chế cốt lõi: các mã đề chỉ khác nhau về
\textit{thứ tự}, nội dung vẫn y hệt; giáo viên thường không đưa được ngân hàng
câu hỏi của riêng mình vào; và kho câu hỏi vẫn là kho tĩnh, chỉ là tĩnh của
người khác.

\textbf{Nhóm 3 --- Dùng trí tuệ nhân tạo sinh thẳng đề.} Đây là hướng đang được
nói đến nhiều nhất hiện nay: đưa yêu cầu cho một mô hình ngôn ngữ lớn và nhận
về một đề hoàn chỉnh. Ưu điểm là nhanh và không cần chuẩn bị gì. Hạn chế cốt
lõi nằm ở chính bản chất của mô hình --- nó dự đoán chuỗi kí tự có xác suất cao
nhất chứ không thực hiện phép toán, nên trong quá trình thử nghiệm tôi thường
xuyên gặp các lỗi sau:

\begin{itemize}[leftmargin=1.2cm, itemsep=2pt]
  \item \textbf{Vượt khung chương trình:} đưa vào kiến thức của lớp trên hoặc
        của chương chưa học, dù yêu cầu đã ghi rõ phạm vi.
  \item \textbf{Tính toán sai:} đáp án không khớp với đề bài, hoặc lời giải
        trình bày trôi chảy nhưng sai ở một bước biến đổi.
  \item \textbf{Không ổn định:} cùng một yêu cầu, chạy hai lần cho hai kết quả
        khác nhau, không kiểm chứng lại được.
  \item \textbf{Sai mức độ:} câu ghi là nhận biết nhưng thực chất là vận dụng,
        làm hỏng ma trận đề.
\end{itemize}

Một đề kiểm tra lấy điểm chính thức không chấp nhận được các rủi ro đó.

\begin{center}
\begin{tabular}{|p{3.2cm}|c|c|c|c|}
\hline
\textbf{Tiêu chí} & \textbf{Ngân hàng tĩnh} & \textbf{Phần mềm trộn đề} &
\textbf{AI sinh đề} & \textbf{Sáng kiến này} \\ \hline
Nội dung do giáo viên quyết định & Có & Không & Không & Có \\ \hline
Sinh được đề mới, khác số liệu & Không & Không & Có & Có \\ \hline
Các mã đề khác nhau thật sự & Không & Không & Có & Có \\ \hline
Đáp án, lời giải chắc chắn đúng & Có & Có & Không & Có \\ \hline
Đúng khung chương trình & Có & Tuỳ nguồn & Không chắc & Có \\ \hline
Đúng ma trận bốn mức độ & Thủ công & Tuỳ nguồn & Không chắc & Tự động \\ \hline
\end{tabular}
\end{center}
\textit{Bảng 3. So sánh sáng kiến với các giải pháp đã có}

Bảng trên cho thấy khoảng trống mà sáng kiến này lấp vào: \textbf{vừa giữ được
quyền quyết định nội dung của giáo viên như ngân hàng tĩnh, vừa sinh được đề
mới như trí tuệ nhân tạo, mà không mang theo rủi ro sai sót của trí tuệ nhân
tạo}. Cách làm để đạt được cả ba điều đó cùng lúc là: câu hỏi do chính giáo
viên viết ra, nhưng viết dưới dạng \textit{chương trình Python sinh câu hỏi}
chứ không phải dưới dạng văn bản cố định. Toàn bộ số liệu, đáp án và lời giải
đều do Python tính ra bằng công thức toán học, nên không thể sai; phạm vi kiến
thức đã bị khoá cứng ngay trong mã định danh của câu hỏi, nên không thể vượt
khung chương trình.

\clearpage

% % ---------------------------------------------------------------
%  Muc II.3 -- Cac giai phap (TRONG TAM)
% ---------------------------------------------------------------
\subsection{Các giải pháp đã thực hiện}

Sáng kiến được thực hiện qua sáu giải pháp nối tiếp nhau, mỗi giải pháp giải
quyết một mắt xích của quy trình: từ cách \textit{đặt tên} câu hỏi, cách
\textit{viết} câu hỏi, cách \textit{quy định} ma trận, cách \textit{chọn} câu,
cách \textit{in} ra đề, đến cách \textit{tổ chức} cho học sinh làm bài.

\subsubsection{Giải pháp 1 --- Chuẩn hoá mã định danh câu hỏi}

Đây là giải pháp nền, quyết định toàn bộ những giải pháp còn lại. Nhận xét
xuất phát rất đơn giản: \textit{muốn máy tính chọn câu hỏi thay cho giáo viên
thì mỗi câu hỏi phải có một địa chỉ mà máy đọc được}. Vì vậy tôi quy định mọi
câu hỏi trong ngân hàng đều mang một mã định danh thống nhất:

\begin{figure}[H]
\centering
\begin{tikzpicture}[font=\small, node distance=0pt]
\node[font=\ttfamily\large] (id) at (0,0) {L10\_C1\_B2\_NB017\_MC\_A};
\node[font=\scriptsize, align=center] (a) at (-4.2,-1.6) {Khối lớp\\ (10/11/12)};
\node[font=\scriptsize, align=center] (b) at (-2.6,-2.4) {Chương};
\node[font=\scriptsize, align=center] (c) at (-1.2,-1.6) {Bài};
\node[font=\scriptsize, align=center] (d) at (0.4,-2.4) {Mức độ\\ NB/TH/VD/VDC};
\node[font=\scriptsize, align=center] (e) at (2.2,-1.6) {Mã năng lực\\ (lấy từ Curriculum)};
\node[font=\scriptsize, align=center] (f) at (3.9,-2.4) {Loại câu\\ MC/TF/SA/TL};
\node[font=\scriptsize, align=center] (g) at (5.2,-1.6) {Phiên bản};
\draw[gray] (-4.05,-0.25) -- (a);
\draw[gray] (-3.05,-0.25) -- (b);
\draw[gray] (-2.05,-0.25) -- (c);
\draw[gray] (-0.75,-0.25) -- (d);
\draw[gray] (0.55,-0.25) -- (e);
\draw[gray] (1.75,-0.25) -- (f);
\draw[gray] (2.75,-0.25) -- (g);
\end{tikzpicture}
\caption{Cấu trúc mã định danh của một câu hỏi trong ngân hàng}
\end{figure}

Bốn quy tắc kèm theo:

\begin{enumerate}[label=(\arabic*), leftmargin=1.2cm, itemsep=3pt]
  \item \textbf{Mã năng lực lấy nguyên văn từ bảng Curriculum}, không được tự
        đặt. Nhờ vậy mỗi câu hỏi luôn gắn với một yêu cầu cần đạt cụ thể của
        Chương trình 2018, và \textit{không thể vượt khung chương trình}: phạm
        vi kiến thức đã bị khoá ngay trong tên của câu hỏi.
  \item \textbf{Câu Đúng/Sai ra theo chương, không theo bài} --- mã rút gọn
        \texttt{L10\_C1\_TF\_A}. Theo quy ước chuyên môn của tôi, một câu
        Đúng/Sai gồm bốn ý ở bốn mức độ nhận thức khác nhau (nhận biết, thông
        hiểu, vận dụng, vận dụng cao), nên khi lên ma trận nó được tính là bốn
        câu chứ không phải một.
  \item \textbf{Mức vận dụng cao mượn mã năng lực của mức vận dụng.} Bảng
        Curriculum chỉ có ba mức nhận biết, thông hiểu, vận dụng; việc một câu
        được ra ở mức vận dụng hay vận dụng cao do bước chọn câu quyết định,
        không sinh thêm mã mới. Quy ước này giúp bảng Curriculum không phình
        ra vô ích.
  \item \textbf{Mã đã phát hành thì không sửa.} Muốn thay nội dung thì tạo
        phiên bản mới (\texttt{\_A} $\to$ \texttt{\_B} $\to$ \texttt{\_C}). Nhờ
        vậy một đề đã in ra từ hai năm trước vẫn tra ngược lại được đúng câu
        hỏi gốc.
\end{enumerate}

Toàn hệ thống --- cơ sở dữ liệu, giao diện, phần mềm điều phối --- đều làm việc
bằng mã định danh này chứ không bằng tên bài tiếng Việt, và \textbf{trí tuệ
nhân tạo tuyệt đối không được phép tự tạo hay tự suy luận ra mã}.

\subsubsection{Giải pháp 2 --- Viết ngân hàng câu hỏi bằng Python}

\textbf{a) Ba tầng dữ liệu}

Ngân hàng không phải một tệp câu hỏi mà là ba tầng dữ liệu nối nhau:

\begin{center}
\begin{tabular}{|p{2.6cm}|p{5.2cm}|p{5.6cm}|}
\hline
\textbf{Tầng} & \textbf{Nội dung} & \textbf{Ví dụ} \\ \hline
Phân phối chương trình & Danh sách bài của từng lớp kèm \textbf{số tiết} và mốc kiểm tra giữa kì, cuối kì &
\texttt{L10\_C1\_B2}, tiết ``23, 25--26'' $\to$ 3 tiết \\ \hline
Curriculum & Các yêu cầu cần đạt của từng bài, chia theo ba mức NB/TH/VD &
\texttt{L10\_C1\_B2\_NB017} \\ \hline
Mapping & Cầu nối: mỗi yêu cầu cần đạt có những hàm sinh câu hỏi nào, thuộc loại nào &
\texttt{L10\_C1\_B2\_NB017\_MC\_A} \\ \hline
\end{tabular}
\end{center}
\textit{Bảng 4. Ba tầng dữ liệu của ngân hàng đề}

\textbf{b) Câu hỏi là một hàm, không phải một đoạn văn bản}

Điểm khác biệt căn bản của sáng kiến nằm ở đây. Mỗi câu hỏi được viết thành một
hàm Python có tên trùng đúng mã định danh của nó. Hàm nhận vào số câu cần sinh
và trả về đoạn \LaTeX\ hoàn chỉnh gồm đề bài, các phương án, đánh dấu phương án
đúng và lời giải chi tiết:

\begin{lstlisting}
def L10_C1_B2_NB017_MC_A_01(socau, socot):
    ket_qua = ""
    for i in range(socau):
        a = random.randint(2, 9)          # sinh so lieu ngau nhien
        b = random.choice([x for x in range(2, 12) if x % a == 0])
        dap_an_dung = b // a              # dap an do PYTHON tinh ra
        ...                               # dung 3 phuong an nhieu
        ket_qua += cau_hoi_latex
    return ket_qua
\end{lstlisting}

Ba hệ quả quan trọng:

\begin{itemize}[leftmargin=1.2cm, itemsep=3pt]
  \item \textbf{Đáp án không thể sai.} Đáp án và lời giải không phải do ai gõ
        vào hay do trí tuệ nhân tạo đoán ra, mà do chính Python tính bằng công
        thức toán học từ bộ số liệu vừa sinh. Mỗi lần chạy là một lần tính
        lại, nên số liệu đổi thì đáp án đổi theo, luôn khớp.
  \item \textbf{Số liệu luôn ``đẹp''.} Phần công sức chuyên môn của giáo viên
        nằm ở chỗ ràng buộc miền giá trị của tham số sao cho mọi biến thể đều
        có nghiệm đẹp, đúng tầm học sinh và không rơi vào trường hợp suy biến.
        Máy không làm thay được việc này.
  \item \textbf{Một mã, nhiều cách ra đề.} Cùng một mã định danh có thể có
        nhiều hàm khác nhau, đánh số \texttt{\_01}, \texttt{\_02},
        \texttt{\_03}\dots\ --- mỗi hàm là một bối cảnh ra đề khác. Khi sinh
        đề, hệ thống chọn ngẫu nhiên một trong các hàm đó. Hậu tố này chỉ tồn
        tại bên trong tệp Python, không lộ ra bất kì nơi nào khác, nên việc bổ
        sung độ phong phú cho ngân hàng không làm thay đổi mã định danh đã
        phát hành.
\end{itemize}

\textbf{c) Kiểm thử tự động toàn bộ ngân hàng}

Ngân hàng càng lớn thì càng khó rà bằng mắt. Tôi viết thêm bộ kiểm thử tự động
gọi lần lượt toàn bộ các hàm sinh câu hỏi và kiểm tra kết quả trả về có đúng
định dạng, có đủ đáp án hay không. Bộ kiểm thử này đã phát hiện một lỗi thật mà
đọc bằng mắt không thấy: một hàm kết thúc bằng \texttt{return} mà quên tên biến
nên trả về rỗng, khiến câu hỏi biến mất khỏi đề trong suốt một thời gian dài.

\chenanh[0.92]{03-ham-sinh-cau-hoi.png}{Một hàm sinh câu hỏi trong ngân hàng Python, kèm phần tính đáp án và lời giải}

\subsubsection{Giải pháp 3 --- Đưa ma trận đề vào bảng cấu hình}

Ma trận đề không được viết lẫn vào mã chương trình mà đặt riêng trong một bảng
cấu hình, để khi Bộ hoặc Sở thay đổi cấu trúc đề thì chỉ cần sửa số trong bảng.
Bảng quy định cho từng loại bài kiểm tra (hệ số 1 --- kiểm tra thường xuyên,
hệ số 2 và 3 --- kiểm tra định kì) số lượng câu của mỗi phần và tỉ lệ bốn mức
độ nhận thức:

\begin{lstlisting}
"HeSo2_HeSo3": {
  "trac_nghiem":     {"so_luong": 12, "ty_le_muc_do": {"NB":0.4,"TH":0.3,"VD":0.2,"VDC":0.1}},
  "dung_sai_cau_lon":{"so_luong":  2, "ty_le_muc_do": {"NB":0.25,"TH":0.25,"VD":0.25,"VDC":0.25}},
  "tra_loi_ngan":    {"so_luong":  4, "ty_le_muc_do": {"NB":0,"TH":0,"VD":0.5,"VDC":0.5}},
  "tu_luan":         {"so_luong":  3, "ty_le_muc_do": {"NB":0,"TH":0,"VD":0.7,"VDC":0.3}}
}
\end{lstlisting}

Một bảng thứ hai quy định thang điểm 10 chia theo phần: trắc nghiệm nhiều lựa
chọn 3 điểm, đúng/sai 2 điểm, trả lời ngắn 2 điểm, tự luận 3 điểm; trong mỗi
phần thì chia đều cho từng câu. Cách tách này cho phép thay đổi trọng số các
phần mà không phải sửa một dòng mã nào.

\ghichu{Nếu tổ chuyên môn của chị thống nhất một ma trận khác (ví dụ hệ số 1
chỉ 6 câu trắc nghiệm), sửa số trong bảng rồi chụp lại ảnh cho khớp với bài
viết.}

\chenanh[0.92]{04-bang-ma-tran.png}{Bảng cấu hình ma trận đề và bảng thang điểm --- thay đổi cấu trúc đề chỉ cần sửa số ở đây}

\subsubsection{Giải pháp 4 --- Thuật toán dựng ma trận chi tiết và chọn câu hỏi}

Đây là phần lõi và cũng là phần được chỉnh sửa nhiều lần nhất. Từ một yêu cầu
của người dùng, hệ thống đi qua sáu bước, \textbf{tất cả đều là chương trình
máy tính, không có trí tuệ nhân tạo tham gia}:

\textbf{Bước 1 --- Xác định phạm vi.} Đọc phân phối chương trình theo lớp.
Nếu là kiểm tra thường xuyên theo chương thì lấy toàn bộ bài của chương đó;
nếu là giữa kì hoặc cuối kì thì lấy từ đầu học kì đến đúng bài có mốc kiểm tra
tương ứng. Bước này trả về danh sách bài kèm \textbf{số tiết} của từng bài.

\textbf{Bước 2 --- Chọn bài cho câu Đúng/Sai trước tiên.} Vì mỗi câu Đúng/Sai
chiếm trọn bốn mức độ nên phải xếp chỗ cho nó trước, nếu không sẽ vỡ ma trận.
Các bài trong phạm vi được sắp theo số tiết giảm dần, lấy từ trên xuống; bốn ý
của một câu Đúng/Sai đều nằm trong cùng một bài.

\textbf{Bước 3 --- Trừ chỉ tiêu ngay tại bài đã chọn.} Một câu Đúng/Sai đã
chiếm 1 nhận biết + 1 thông hiểu + 1 vận dụng + 1 vận dụng cao, nên phải trừ
đúng bốn ý đó ra khỏi chỉ tiêu của chính bài đó --- \textit{``đơn vị kiến thức
đó chia ra được 3 câu thì phải tính 1 ở đúng sai, chỉ còn 2''}. Trừ ở cấp bài
đã tự động làm tổng giảm đúng bằng số câu Đúng/Sai, tuyệt đối không trừ thêm
lần thứ hai ở cấp tổng.

\textbf{Bước 4 --- Chia câu về từng bài theo tỉ lệ số tiết.} Các câu mức nhận
biết và thông hiểu được chia về từng bài \textit{theo tỉ lệ số tiết của bài đó
trong phân phối chương trình} --- bài dạy nhiều tiết thì nhiều câu, đúng
nguyên tắc trọng số của ma trận đề. Phần dư sau khi làm tròn xuống được rải
lần lượt cho bài nhiều tiết nhất trước.

\textbf{Bước 5 --- Rải mức vận dụng và vận dụng cao.} Chọn vận dụng cao trước,
\textbf{mỗi bài nhiều nhất một câu vận dụng cao} để không dồn hết câu khó vào
một đơn vị kiến thức. Sau đó mới rải câu vận dụng, ưu tiên những bài chưa có
câu vận dụng cao. Tỉ lệ vận dụng so với vận dụng cao lấy đúng từ bảng cấu hình
ở Giải pháp 3, không ép cứng trong mã.

\textbf{Bước 6 --- Chọn mã câu hỏi cuối cùng.} Đến đây ma trận đã nói rõ ``cần
một câu mức thông hiểu, thuộc yêu cầu cần đạt số 014 của bài 1, loại trắc
nghiệm nhiều lựa chọn''. Bước chọn câu tra bảng Mapping để tìm mã hàm sinh
tương ứng, bỏ qua những mã đã dùng trong đề hiện tại, và xoay vòng giữa các
phiên bản A, B, C để đề nào cũng khác đề nào.

Sau đó câu hỏi sinh ra còn qua một bước kiểm tra tự động: đúng chương, đúng
bài, đúng mức độ, đúng loại câu, không trùng câu trong cùng một đề.

\textbf{Chế độ nháp.} Nếu ngân hàng chưa đủ câu cho một ô nào đó của ma trận,
hệ thống không dừng cả đề mà chèn dòng ``[THIẾU CÂU HỎI]'' vào đúng vị trí.
Đây là lựa chọn có chủ ý: dòng đó vừa cho giáo viên một đề dùng tạm được, vừa
là chỉ dấu rõ ràng cho biết cần bổ sung hàm sinh câu hỏi nào tiếp theo.

\chenanh[0.92]{05-ma-tran-de-sinh-ra.png}{Ma trận chi tiết do hệ thống dựng ra cho một đề kiểm tra định kì}

\clearpage
\subsubsection{Giải pháp 5 --- Sinh đề, lời giải và nhiều mã đề bằng \LaTeX}

\textbf{a) Đề và lời giải từ cùng một nguồn}

Các câu hỏi đã chọn được ghép thành một tệp \LaTeX\ theo mẫu đề thi, rồi biên
dịch ra PDF. Bản lời giải không được sinh lại lần nữa mà lấy chính tệp
\LaTeX\ của đề vừa in, chỉ đổi chế độ hiển thị từ ``đề thi'' sang ``lời giải''.
Nhờ vậy \textbf{lời giải chắc chắn ứng đúng với đề học sinh đang cầm trên tay},
không bao giờ xảy ra tình trạng lệch câu giữa đề và đáp án --- một lỗi rất
thường gặp khi làm thủ công.

Đáp án cũng được trích ngay tại thời điểm sinh đề, đọc trực tiếp từ các dấu
hiệu trong mã \LaTeX\ (phương án được đánh dấu đúng, đáp số của câu trả lời
ngắn, phần lời giải), rồi lưu lại thành một tệp đáp án đi kèm. Khi chấm bài,
hệ thống so với tệp đáp án này, nên kết quả chấm khớp 100\% với đề đã phát.

\textbf{b) Nhiều mã đề cùng cấu trúc, chỉ khác số liệu}

Đây là chức năng được giáo viên đánh giá cao nhất khi dùng thử. Vì mỗi câu hỏi
là một hàm sinh, nên chỉ cần gọi lại hàm đó là có ngay một câu mới:
\textbf{cùng bài, cùng yêu cầu cần đạt, cùng mức độ nhận thức, cùng dạng toán
--- chỉ khác số liệu}. Từ đó hệ thống in ra bao nhiêu mã đề cũng được, và các
mã đề ấy tương đương nhau tuyệt đối về nội dung lẫn độ khó.

Ý nghĩa thực tế với giáo viên: \textit{lấy ra là dùng kiểm tra được ngay, không
phải ngồi tự thay số để tạo đề mới, cũng không phải giải lại để kiểm tra đáp
án}. Ý nghĩa với học sinh: hai em ngồi cạnh nhau nhận hai đề khác nhau thật
sự, nhưng không em nào thiệt hơn em nào.

Cũng từ cơ chế đó, trang kết quả sau khi làm bài có thêm nút ``Làm đề khác cùng
cấu trúc''. Học sinh bấm một cái là có ngay một đề mới y hệt về cấu trúc để
luyện lại phần mình vừa làm sai --- và đường đi này \textbf{không gọi trí tuệ
nhân tạo}, vì tham số của đề cũ đã được lưu sẵn, nên vừa nhanh vừa không tốn
chi phí.

\chenanh[0.92]{06-hai-ma-de.png}{Hai mã đề của cùng một cấu trúc: cùng dạng toán, cùng mức độ, chỉ khác số liệu}

\chenanh[0.92]{07-de-va-loi-giai-pdf.png}{Tệp PDF đề thi và tệp PDF lời giải tương ứng do hệ thống sinh ra}

\subsubsection{Giải pháp 6 --- Tổ chức làm bài, chấm điểm và theo dõi kết quả}

\textbf{a) Làm bài trực tuyến và chấm ngay.} Học sinh mở đề trên web, làm các
phần trắc nghiệm nhiều lựa chọn, đúng/sai và trả lời ngắn ngay trên trình
duyệt. Bấm nộp là có điểm ngay cùng với đáp án đúng của từng câu. Việc chấm do
chương trình đối chiếu với tệp đáp án đã lưu, \textbf{không gọi trí tuệ nhân
tạo}, nên vừa chính xác vừa không phát sinh chi phí dù bao nhiêu học sinh cùng
làm bài.

\textbf{b) Thang điểm 10 chia theo phần.} Điểm được tính theo đúng bảng trọng
số ở Giải pháp 3: trắc nghiệm nhiều lựa chọn 3 điểm, đúng/sai 2 điểm, trả lời
ngắn 2 điểm, tự luận 3 điểm, mỗi phần chia đều cho các câu trong phần đó. Hiện
tại phần tự luận 3 điểm học sinh làm ra giấy nộp cho giáo viên chấm, nên màn
hình kết quả ghi rõ điểm tự động là 7 điểm và chức năng chấm tự luận đang được
xây dựng --- ghi rõ như vậy để học sinh không hiểu nhầm là mình bị mất điểm.

\textbf{c) Giao đề và trả điểm qua Google Classroom.} Sau khi giáo viên kết
nối một lần tài khoản Google Classroom của mình, hệ thống tự đăng đề, lời giải
và điểm của từng học sinh về đúng lớp học thật trên Classroom. Học sinh mở
Classroom quen thuộc là thấy bài của mình, không phải học thêm một phần mềm
mới.

\textbf{d) Trang thống kê cho giáo viên.} Giáo viên xem được điểm trung bình
của lớp, danh sách điểm từng học sinh và --- quan trọng hơn cả --- những
chương, bài mà học sinh sai nhiều nhất. Đây chính là thông tin để điều chỉnh
việc dạy: bài nào cả lớp sai nhiều thì cần dạy lại, em nào sai lệch hẳn so với
lớp thì cần kèm riêng.

\chenanh[0.92]{15-lam-bai-nop-bai.png}{Màn hình làm bài trực tuyến và màn hình kết quả chấm ngay sau khi nộp}

\subsubsection{Công cụ và hạ tầng đã sử dụng}

Sáng kiến này là một phần mềm chạy thật trên Internet, nên tôi xin nêu rõ toàn
bộ công cụ đã dùng để đồng nghiệp muốn làm theo biết cần chuẩn bị những gì.

\begin{center}
\begin{longtable}{|p{3.6cm}|p{4.6cm}|p{6.2cm}|}
\hline
\textbf{Thành phần} & \textbf{Công cụ đã dùng} & \textbf{Dùng để làm gì} \\ \hline
\endhead
Ngôn ngữ lập trình & Python 3 & Viết hàm sinh câu hỏi và toàn bộ phần xử lí \\ \hline
Máy chủ web & FastAPI + Jinja2 & Cung cấp giao diện và các đường dẫn của trang web \\ \hline
Cơ sở dữ liệu, đăng nhập & Supabase (PostgreSQL) & Lưu tài khoản, lịch sử hội thoại, đề đã sinh, bài làm, điểm \\ \hline
Điều phối luồng & n8n (bản tự cài trên máy chủ riêng) & Nối các bước với nhau và là nơi duy nhất gọi tới trí tuệ nhân tạo \\ \hline
Mô hình ngôn ngữ & Gọi qua OpenRouter (Qwen, Gemini) & Chỉ để hiểu yêu cầu tiếng Việt của người dùng \\ \hline
Trình bày đề & \LaTeX, gói \texttt{ex\_test}, biên dịch bằng \texttt{pdflatex} & Sinh tệp PDF đề thi và lời giải đúng chuẩn văn bản đề \\ \hline
Máy chủ & Máy ảo (VPS) chạy Ubuntu, nginx, systemd & Nơi phần mềm chạy 24/24 \\ \hline
Tên miền & \texttt{nganhangdechv.tech} (tự mua) & Địa chỉ để giáo viên và học sinh truy cập \\ \hline
Tích hợp ngoài & Google Classroom API, Google Drive API & Giao đề, đáp án và trả điểm về lớp học thật \\ \hline
Quản lí mã nguồn & Git và GitHub & Lưu lịch sử sửa đổi, đồng bộ giữa máy cá nhân và máy chủ \\ \hline
Soạn thảo & Visual Studio Code & Nơi viết toàn bộ mã nguồn \\ \hline
\end{longtable}
\end{center}
\textit{Bảng 5. Công cụ và hạ tầng đã sử dụng}

Trong đó, hai thành phần cần giải thích thêm:

\textbf{n8n} là một phần mềm nối các bước xử lí thành một luồng, thao tác bằng
cách kéo thả các khối trên màn hình thay vì viết mã. Tôi cài n8n trên máy chủ
của mình (bản tự cài, không phải bản thuê theo tháng). Trong hệ thống này, n8n
giữ đúng một vai trò: nhận yêu cầu tiếng Việt, chuyển cho mô hình ngôn ngữ hiểu,
rồi gọi sang phần xử lí bằng Python. Mọi việc có quy tắc rõ ràng đều đã được
chuyển hết về Python, đúng nguyên tắc ``ưu tiên mã nguồn hơn trí tuệ nhân tạo''.

\textbf{Máy ảo (VPS)} là một máy tính thuê trên Internet, chạy liên tục, để
phần mềm luôn sẵn sàng khi học sinh truy cập bất kể giờ nào. Cùng với tên miền,
đây là hai khoản chi phí duy nhất của sáng kiến, do tôi tự chi trả.

\ghichu{Nếu muốn ghi rõ chi phí, điền số tiền thật mỗi năm cho máy ảo và tên
miền vào đoạn trên --- con số cụ thể sẽ làm phần ``hiệu quả kinh tế'' ở Mục
II.5 thuyết phục hơn nhiều.}

\textbf{Mã nguồn công khai.} Toàn bộ mã nguồn của hệ thống, kèm hơn 20 tài liệu kĩ
thuật mô tả kiến trúc, chuẩn mã định danh, thuật toán dựng ma trận và nhật kí
toàn bộ các lần cải tiến, được công khai tại:

\begin{center}
\setlength{\fboxsep}{8pt}
\fbox{\begin{minipage}{0.85\linewidth}\centering
\textbf{\url{https://github.com/lantuan/NganHangDe}}
\end{minipage}}
\end{center}

Tôi công khai mã nguồn vì hai lí do. Thứ nhất, để hội đồng và đồng nghiệp kiểm
chứng được rằng sản phẩm là có thật và do tôi làm, chứ không chỉ là mô tả trên
giấy --- riêng nhật kí sửa đổi trong mã nguồn đã ghi lại từng lần cải tiến ở
Bảng 6 dưới đây. Thứ hai, để giáo viên nào muốn làm theo có thể lấy về dùng
lại, không phải bắt đầu từ con số không.

\subsubsection{Quá trình cải tiến}

Sáng kiến không hoàn chỉnh ngay từ bản đầu tiên. Toàn bộ quá trình được ghi lại
trong sổ tay kĩ thuật của hệ thống theo từng phiên bản; dưới đây là những lần
cải tiến đáng kể nhất, đi kèm lí do --- vì chính những lần sửa này mới là nội
dung có giá trị chia sẻ với đồng nghiệp.

\begin{center}
\begin{longtable}{|p{1.4cm}|p{5.6cm}|p{6.4cm}|}
\hline
\textbf{Lần} & \textbf{Vấn đề gặp phải} & \textbf{Cách khắc phục} \\ \hline
\endhead
1 & Bản thiết kế ban đầu định dùng nhiều trí tuệ nhân tạo cho từng khâu (phân
tích yêu cầu, dựng ma trận, chọn câu) --- kết quả không ổn định, mỗi lần chạy
một khác &
Rút xuống chỉ còn ba chỗ dùng trí tuệ nhân tạo, toàn bộ khâu còn lại chuyển
sang chương trình máy tính. Đây là lần cải tiến quan trọng nhất, đặt ra nguyên
tắc xuyên suốt của hệ thống \\ \hline

2 & Đáp án phải nhập tay riêng nên hay lệch với đề &
Phát hiện đáp án đã nằm sẵn trong mã \LaTeX\ của câu hỏi; viết chương trình tự
trích ra ngay lúc sinh đề. Việc trích không dùng được cách tìm kiếm thông
thường vì công thức toán có ngoặc lồng nhau, phải viết bộ đếm ngoặc riêng \\ \hline

3 & Lời giải bị lọt vào đề của học sinh &
Bổ sung mốc chặn khi cắt phần đề bài, đồng thời lược bỏ các lệnh \LaTeX\ mà
trình duyệt không hiển thị được \\ \hline

4 & Mọi câu Đúng/Sai in ra PDF bị dính liền thành một đoạn văn &
Truy ra một lệnh \LaTeX\ được sinh ra nhưng không hề tồn tại trong gói lệnh
đang dùng; sửa ở thư viện sinh câu hỏi để luôn dùng lệnh đúng \\ \hline

5 & Toàn bộ câu trả lời ngắn bị sinh nhầm thành câu trắc nghiệm ở mọi đề đã
từng tạo &
Nguyên nhân là chỗ gọi hàm luôn truyền cùng một giá trị cho tham số dạng câu;
sửa lại nhánh gọi và bổ sung giá trị mặc định cho 29 hàm sinh \\ \hline

6 & Một hàm sinh câu hỏi trả về rỗng, câu hỏi biến mất khỏi đề mà không ai biết &
Viết bộ kiểm thử tự động chạy qua toàn bộ các hàm trong ngân hàng; lỗi bị phát
hiện ngay và từ đó mọi lần bổ sung câu hỏi mới đều phải qua bộ kiểm thử này \\ \hline

7 & Số câu do người dùng yêu cầu bị bỏ mất trên đường truyền, đề luôn ra theo
số mặc định &
Rà lại toàn bộ đường đi của dữ liệu, phát hiện một bước trung gian gán cứng giá
trị rỗng đè lên yêu cầu thật \\ \hline

8 & Trí tuệ nhân tạo hiểu ``chương 1'' thành ``lớp 1'' &
Chặn sớm ngay đầu vào: lớp không thuộc 10, 11, 12 thì báo lỗi rõ ràng; sau đó
đổi hẳn thành hỏi lại người dùng thay vì tự đoán \\ \hline

9 & Câu ``cho tôi thêm đề nữa'' bị hiểu là ngoài phạm vi &
Thay vì chỉnh trí tuệ nhân tạo cho hiểu, làm hẳn một đường đi không qua trí tuệ
nhân tạo: nút ``Làm đề khác cùng cấu trúc'' dùng lại tham số của đề cũ \\ \hline

10 & Số câu chia đều cho các bài, không đúng trọng số &
Chuyển sang chia theo \textbf{tỉ lệ số tiết} của từng bài trong phân phối
chương trình \\ \hline

11 & Câu Đúng/Sai bị trừ chỉ tiêu hai lần, làm hụt câu ở các phần khác &
Quy định trừ đúng một lần và trừ ngay tại bài mà câu Đúng/Sai đã chọn \\ \hline

12 & Các câu vận dụng cao dồn hết vào một đơn vị kiến thức &
Quy định mỗi bài nhiều nhất một câu vận dụng cao, chọn vận dụng cao trước rồi
mới rải vận dụng \\ \hline

13 & Học sinh mất phiên đăng nhập giữa chừng &
Bổ sung cơ chế tự làm mới phiên khi phiên hết hạn \\ \hline

14 & Đề và đáp án không hiện trong tài khoản Classroom của học sinh &
Ba nguyên nhân nối nhau: giao diện nuốt mất thông báo lỗi; chưa bật một dịch vụ
cần thiết của Google; và xin nhầm quyền (quyền đăng \textit{bài tập} khác quyền
đăng \textit{tài liệu}). Sửa lần lượt cả ba, đồng thời cho hiện lí do thất bại
ngay trên màn hình thay vì im lặng \\ \hline
\end{longtable}
\end{center}
\textit{Bảng 6. Những lần cải tiến đáng kể trong quá trình thực hiện}

Nhìn lại, hầu hết các lỗi đều thuộc một trong hai nhóm: \textit{giao cho trí
tuệ nhân tạo việc lẽ ra phải do chương trình làm}, và \textit{không có cơ chế
tự kiểm tra nên lỗi âm thầm tồn tại rất lâu}. Hai bài học rút ra tương ứng là
nguyên tắc phân vai ở Giải pháp 1 và bộ kiểm thử tự động ở Giải pháp 2.

\subsubsection{Vai trò của trợ lí trí tuệ nhân tạo trong quá trình thực hiện}

Tôi xin nói rõ và không né tránh: trong quá trình xây dựng hệ thống này, tôi có
sử dụng trợ lí trí tuệ nhân tạo (Claude của Anthropic) như một cộng sự kĩ
thuật. Cụ thể, trợ lí giúp tôi viết và sửa mã nguồn theo yêu cầu, dò tìm
nguyên nhân của các lỗi nêu ở bảng trên, và ghi lại sổ tay kĩ thuật của dự án.

Tuy nhiên, phần thuộc về chuyên môn sư phạm vẫn hoàn toàn do tôi quyết định và
chịu trách nhiệm: quy ước mã định danh câu hỏi, quy ước một câu Đúng/Sai gồm
bốn ý ở bốn mức độ, quy tắc trừ chỉ tiêu tại chính bài đã chọn, nguyên tắc chia
câu theo tỉ lệ số tiết, quy định mỗi bài nhiều nhất một câu vận dụng cao, ma
trận và thang điểm của từng loại bài kiểm tra, cùng toàn bộ nội dung toán học
và ràng buộc số liệu của từng hàm sinh câu hỏi. Nhiều lần trợ lí đề xuất cách
làm khác, tôi phải yêu cầu làm lại cho đúng quy ước sư phạm mà tôi đã đặt ra.

Chính trải nghiệm đó củng cố thêm quan điểm nêu ở Mục II.1: trí tuệ nhân tạo là
một công cụ rất mạnh để hiện thực hoá ý tưởng của giáo viên, nhưng không thay
được người giáo viên trong việc quyết định \textit{dạy gì, kiểm tra gì và kiểm
tra đến mức nào}. Với các đồng nghiệp muốn làm theo hướng này, tôi cho rằng
điều kiện cần không phải là biết lập trình giỏi, mà là \textbf{có quy ước
chuyên môn rõ ràng trước khi ngồi vào máy} --- vì đó chính là thứ mà công cụ
không tự nghĩ ra được.

\clearpage

% \input{noidung/06_hieu_qua}
% \input{noidung/07_ket_luan}
% \input{noidung/08_tai_lieu_phu_luc}

\end{document}
